\documentclass[11pt]{article}

\usepackage{amsmath,amssymb,graphicx}
\usepackage[margin=1in]{geometry}
\usepackage{booktabs}

\title{Effective Binding Existence Tests on a Frozen Kernel-Induced Cochain Operator Architecture}
\author{Tomislav Rupi\'c \\ Independent Researcher}
\date{March 2026}

\begin{document}

\maketitle

\begin{abstract}
Stages 10--13 established a frozen pre-geometric morphology program on a kernel-induced cochain operator architecture, including single-packet transport classes, collective interaction morphology, coarse-grained envelope structure, and coupled resolution--interaction diagnostics. The strongest collective representatives from that program support metastable composite envelopes and refinement-driven sharpening of coarse basin structure, but they do not yet yield scale-robust composite interaction labels. Stage 14 therefore shifts from atlas construction to a focused existence test: can one minimal additional structural ingredient produce measurable improvement in composite persistence relative to the frozen Stage 13 baseline? Three narrow branches are tested on the same three representatives: weak local envelope feedback, weak phase-lock retention, and weak envelope-mediated coupling. Each branch is run as a 9-run null-versus-perturbed pilot. In all three branches, every run remains in the same interaction and coarse morphology class as its null control and no run satisfies the promotion criterion for refinement follow-up. For Branch A, a direct term-magnitude probe shows that the added feedback survives projection almost intact while remaining small relative to the projected base dynamics, so its null result is a true negative at the tested strengths. The combined Stage 14 read is therefore bounded and simple: under the tested weak structural extensions, no effective binding channel appears, and collective composite persistence remains a superposition-morphology phenomenon rather than a stabilized interaction mechanism.
\end{abstract}

\section{Introduction}

Stages 10--13 established a layered diagnostic arc on a frozen kernel-induced cochain operator architecture. Stage 10 defined the morphology instrument for single-packet transport and recorded its baseline, resilience, and transition structure. Stage 11 extended the same packet family to collective interaction morphology and showed that moderate-separation configurations occupy a transient binding corridor while tight overlap can produce metastable composite envelopes. Stage 12 coarse-grained these collective configurations and showed that compact clusters can sharpen into basin-dominated envelope structure. Stage 13 then coupled resolution changes directly to interaction and coarse-graining scales and found that coarse basin organization strengthens under refinement while composite interaction labels remain fragile, especially under width coupling.

That progression leaves a precise unresolved question. If the frozen linear architecture supports metastable composites and coarse envelope organization, can a minimal local structural extension stabilize those composites into a more persistent interaction channel? Stage 14 is the first focused existence test of that possibility.

The design principle is intentionally narrow. Stage 14 is not another atlas sweep. It does not introduce a broad new parameter space, and it does not attempt continuum, force, particle, or geometric interpretation. Instead, it tests a small set of weak structural branches against the strongest frozen Stage 13 representatives and asks whether any of them produce a measurable gain in collective composite persistence.

\section{Frozen Baseline and Stage 14 Representatives}

All Stage 14 runs preserve the frozen operator architecture. The projected transverse sector, periodic boundary condition, Gaussian kernel construction, and packet family inherited from the strongest coherent transverse seed remain unchanged. The common packet family is characterized by bandwidth $w=0.1$, cosine carrier, central wave number $k=1.0$, and amplitude scale $0.5$.

Three Stage 13 representatives are used:
\begin{itemize}
\item tight clustered pair,
\item asymmetric density cluster,
\item counter-propagating corridor.
\end{itemize}

The first two represent the strongest composite corridor available after Stages 11--13. The third is a diffuse control configuration that remains transiently bound but does not enter the same composite envelope class.

The frozen null-control behavior for these representatives is summarized in Table~\ref{tab:baseline}.

\begin{table}[h]
\centering
\small
\begin{tabular}{p{0.24\textwidth}p{0.22\textwidth}p{0.20\textwidth}cc}
\toprule
Representative & Interaction label & Coarse label & Lifetime & Binding score \\
\midrule
Tight clustered pair & Metastable composite & Diffuse coarse field & 1.4764 & 1.0000 \\
Asymmetric density cluster & Metastable composite & Diffuse coarse field & 1.4764 & 1.0000 \\
Counter-propagating corridor & Transient binding & Diffuse coarse field & 0.9842 & 0.6667 \\
\bottomrule
\end{tabular}
\caption{Frozen Stage 13 baseline representatives reused in Stage 14. Lifetime denotes composite lifetime in the Stage 11 collective observable grammar. Binding score denotes the Stage 14 binding persistence score.}
\label{tab:baseline}
\end{table}

\section{Stage 14 Design}

Each Stage 14 branch is evaluated against a null control with identical initial conditions and resolution. The pilot format is the same in every branch:
\begin{itemize}
\item three representatives,
\item one null control,
\item one very weak branch strength,
\item one weak branch strength,
\item total of $9$ runs per branch.
\end{itemize}

The same Stage 11 and Stage 12 measurement layers are reused wherever possible:
\begin{itemize}
\item pairwise centroid distance,
\item overlap persistence,
\item composite lifetime,
\item interaction morphology label,
\item dominant basin area fraction,
\item basin count stability,
\item coarse morphology label.
\end{itemize}

Stage 14 adds a small existence-test decision gate. A run is promoted to a refinement follow-up only if it improves composite lifetime, binding persistence, and coarse persistence beyond fixed thresholds relative to its null control while not degrading the interaction label. No branch satisfied that gate in the present study.

The three tested branches are:
\begin{description}
\item[Branch A.] Weak local envelope feedback acting through a local envelope-dependent modifier in the transverse dynamics.
\item[Branch B.] Weak phase-lock retention acting through a local alignment-memory profile in overlap regions.
\item[Branch C.] Weak envelope-mediated coupling acting through component-wise coupling to the smoothed envelope of the other packets.
\end{description}

The tested strengths are deliberately small. Branch A uses coefficients $\varepsilon = 0.015$ and $0.03$. Branches B and C use coefficients $\varepsilon = 0.03$ and $0.06$. The aim is not to force visible change, but to test whether a minimal extension is already sufficient to generate an existence signal.

\section{Branch A: Local Envelope Feedback}

Branch A compares null control against very weak and weak local envelope feedback on all three representatives. The outcome is maximally simple: all nine runs remain identical to their null controls at the level of the recorded Stage 14 metrics. Every run retains the same interaction label, the same coarse label, the same composite lifetime, and the same binding persistence score as the corresponding null control. The branch-level robustness count is therefore
\[
\texttt{no composite gain} = 9.
\]
No run is promoted to a refinement follow-up.

A null result of this cleanliness raises an obvious numerical question: is the added term truly present, or is it projected away? To resolve that ambiguity, a direct term-magnitude probe was run on the tight clustered pair at weak feedback and resolution $n=12$. The probe measured the raw feedback norm, the projected feedback norm, and the ratio of projected feedback norm to projected base-force norm over sampled time slices. The result is decisive:
\begin{itemize}
\item mean projection-survival ratio: $0.9750$,
\item survival range: $0.9747$ to $0.9753$,
\item mean projected-to-base-force ratio: $0.0219$,
\item maximum projected-to-base-force ratio: $0.0226$.
\end{itemize}

So the projector does not annihilate the branch. Roughly $97.5\%$ of the weak feedback norm survives projection, and the surviving term remains at the level of about $2.2\%$ of the projected base-force norm. The correct read is therefore not ``projection-suppressed null.'' It is a true negative at the tested coefficients: a weak local envelope feedback of that scale does not measurably improve collective composite persistence.

\section{Branch B: Phase-Lock Retention}

Branch B replaces the local envelope feedback with a weak memory-like retention of phase alignment in overlap regions. The same three representatives and null-control format are used, with coefficients $\varepsilon = 0.03$ and $0.06$.

Again the result is clean. All nine runs remain identical to their null controls at the level of the recorded metrics. The two composite representatives remain metastable composite regimes with the same lifetime and binding score as the null controls, while the counter-propagating corridor remains a transient binding regime with the same lifetime and binding score as its null control. The coarse label stays diffuse coarse field in all nine runs. No run is promoted to a refinement follow-up, and the robustness count is again
\[
\texttt{no composite gain} = 9.
\]

At the tested strengths, weak phase-lock retention therefore produces no measurable improvement in the Stage 14 existence-test instrument.

\section{Branch C: Envelope-Mediated Coupling}

Branch C keeps the frozen base operator but adds a weak auxiliary coupling channel computed from the smoothed envelope of the other packets. The same three representatives and null-control format are used, with coefficients $\varepsilon = 0.03$ and $0.06$.

The outcome matches Branch B. All nine runs are unchanged at the level of the recorded metrics. The two composite representatives remain metastable composite regimes, the corridor control remains transient binding, the coarse label remains diffuse coarse field throughout, and no run meets the promotion criterion for refinement follow-up. The robustness count is again
\[
\texttt{no composite gain} = 9.
\]

At the tested strengths, weak envelope-mediated coupling does not generate a measurable persistence gain either.

\section{Combined Stage 14 Read}

Across the three branches, Stage 14 contains $27$ base comparisons:
\begin{itemize}
\item $9$ for Branch A,
\item $9$ for Branch B,
\item $9$ for Branch C.
\end{itemize}

The combined result is unusually clean:
\begin{itemize}
\item no run improves composite lifetime relative to null,
\item no run improves binding persistence relative to null,
\item no run improves coarse persistence relative to null,
\item no run changes interaction or coarse morphology class,
\item no run is promoted to a refinement follow-up,
\item every run is labeled \texttt{no composite gain}.
\end{itemize}

For the tight clustered pair and asymmetric density cluster, every tested branch preserves the same metastable composite label and the same lifetime as the frozen baseline. For the counter-propagating corridor, every tested branch preserves the same transient binding label and the same shorter lifetime as the frozen baseline. In this sense, all three tested structural extensions are dynamically silent at the level measured by the existing collective and coarse instruments.

The strongest formal statement Stage 14 supports is therefore limited but useful: under the tested weak structural extensions, no effective binding existence signal appears.

\section{Interpretation Boundary}

Stage 14 remains a diagnostic study. It does not support:
\begin{itemize}
\item particle language,
\item mass or force language,
\item gravitational interpretation,
\item true physical bound-state claims,
\item spacetime or continuum claims.
\end{itemize}

The only bounded claims supported by the present data are:
\begin{itemize}
\item the frozen Stage 13 representatives remain stable enough to serve as Stage 14 null controls,
\item weak local envelope feedback is a true negative at the tested coefficients,
\item weak phase-lock retention shows no measurable effect at the tested coefficients,
\item weak envelope-mediated coupling shows no measurable effect at the tested coefficients,
\item the collective composite persistence observed through Stage 13 behaves, at this level, as a superposition-morphology phenomenon rather than a stabilized interaction mechanism.
\end{itemize}

\section{Conclusion}

Stage 14 is the first focused existence test in the program. Instead of enlarging the atlas, it asks a narrower question: can one minimal additional structural ingredient stabilize the strongest Stage 13 composite representatives into a more persistent interaction channel? Three weak branches were tested against frozen null controls: local envelope feedback, phase-lock retention, and envelope-mediated coupling.

The answer is negative at the tested strengths. Branch A is a true negative, with direct term-magnitude diagnostics showing that the added term survives projection but remains too weak to change the observed morphology. Branches B and C are also null within the present measurement instrument: no metric, label, or promotion gate changes across their 9-run pilots. The combined Stage 14 read is therefore both conservative and informative. The collective structures found in Stages 11--13 remain real as morphology, but they do not yet become stabilized effective composites under the weakest tested structural extensions.

This closes Stage 14 as a falsification-oriented layer. The next useful direction, if the program continues, is not to repeat weak coefficient sweeps. It is to test a different structural hypothesis altogether, such as delayed feedback, nonlocal memory, or scale-adaptive coupling, while preserving the same null-control discipline.

\end{document}
